As of March~22,~2026, the team has moved from specification into \textbf{working vertical slices}: an Expo client, FastAPI API, and PostgreSQL database support authenticated use of the main tabs and core flows described in the proposal.

\subsection{Completed milestones and components}
\begin{itemize}[leftmargin=*]
  \item \textbf{Proposal and design specification:} Scope, problem statement, five-tab UI (Home, Calendar, Add, Grocery, Profile), wireframe, and timeline documented.
  \item \textbf{Backend foundation:} Application entry (\texttt{main.py}, \texttt{run.py}), models for users, recipes (ingredients, steps), meal plans, grocery lists/items, and versioned migrations.
  \item \textbf{Authentication:} Signup and login on the client; JWT attached via an Axios interceptor (\texttt{frontend/services/api.ts}); protected creation of recipes and user-scoped grocery access patterns begun.
  \item \textbf{Recipe library (Home):} Scrollable list of recipe cards (\texttt{RecipeList}, \texttt{RecipeCard}) opening recipe detail by id.
  \item \textbf{Recipe detail:} Hero image when \texttt{image\_url} is set, summary stats, ingredients, and steps; \textbf{Begin Cooking} navigates to step-by-step mode (\texttt{app/recipe/[id].tsx}, \texttt{app/cooking/[id].tsx}).
  \item \textbf{Cooking mode:} One step at a time with previous/next, progress bar, optional per-step timer start when \texttt{timer\_seconds} is defined; \texttt{TimerDock} shows active timers sorted by least time remaining with play/pause and dismiss (see \texttt{TimerDock.tsx} and \texttt{useTimerStore.ts} under \texttt{frontend/}).
  \item \textbf{Add recipe:} Manual form with dynamic ingredient and step rows and optional timer fields per step (\texttt{app/add/manual.tsx}); paste flow opens a modal, calls \texttt{POST /recipes/parse}, and pre-fills the manual screen for verification (\texttt{app/(tabs)/add.tsx}). Image scan is a disabled \emph{Coming Soon} control.
  \item \textbf{Meal planner:} API range query and create plan; UI lists entries for a selected date and supports quick add of custom food name and calories (\texttt{app/(tabs)/calendar.tsx}). The \texttt{MealPlan} model supports \texttt{recipe\_id}; the current UI emphasizes quick-add over picking from the library.
  \item \textbf{Grocery list:} Per-user list, items grouped by category in a sectioned list, checkbox state persisted through the API, manual add modal (\texttt{app/(tabs)/grocery.tsx}, \texttt{endpoints/grocery.py}).
  \item \textbf{Profile:} Displays name, email, and calorie goal from \texttt{GET /users/me} (\texttt{app/(tabs)/profile.tsx}). Schema fields such as \texttt{dietary\_restrictions} exist for future editing and filtering.
\end{itemize}

\subsection{Gaps relative to the proposed UI}
\begin{itemize}[leftmargin=*]
  \item \textbf{Home:} Search and filter icons in the app bar are not wired to queries or client-side filters.
  \item \textbf{Cooking mode:} Per-step mise en place is not modeled (\texttt{Ingredient} is not linked to \texttt{Step}); the screen emphasizes instruction text and timers only.
  \item \textbf{Calendar:} Proposal called for a monthly grid with indicators and a bottom panel to add meals from the library; the current UI is a date-centric list with quick-add only.
  \item \textbf{Grocery:} No server or client pipeline yet to aggregate ingredients from planned recipes (deduplication, scaling, categories).
  \item \textbf{AI parsing:} Deterministic mock output, not a production LLM or vision pipeline---sufficient for flow testing only.
  \item \textbf{Profile:} No in-app editing of restrictions, goals, or social features despite partial schema support.
\end{itemize}

This aligns broadly with the proposal's phasing: Weeks~1--4 emphasized core prototype work (library, cooking mode with timer dock, planner groundwork), while grocery automation, richer input, and profile/settings were slated for subsequent weeks---though basic Add and Grocery screens were started early.
